\section{Sissejuhatus} \label{sissejuhatus}

Bakalaureusetöö kirjutamine on iga informaatika tudengi õpingute lõpus tehtav mahukas iseseisev töö, mille puhul hinnatakse nii teaduslikku panust kui ka töö vormistamist ja keelelist kvaliteeti. Tartu Ülikooli arvutiteaduse instituudi nõuetes moodustab vormistamine ligi 25~\% lõputöö koondhindest~\cite{ut_juhend_2024}, kusjuures suurt rõhku pannakse selgele struktuurile, akadeemilisele stiilile, terminoloogia järjepidevusele ja korrektsele viitamisele. Eestikeelses informaatika lõputöös on neid nõudeid eriti keeruline täita, sest valdkonna terminoloogia on suuresti arenenud välja inglise keeles ning paljudele mõistetele puudub väljakujunenud eestikeelne vaste või on neid mitu.

Suurte keelemudelite (\emph{ingl.\ large language model}, edaspidi LLM) viimaste aastate areng on toonud kaasa loomuliku keele töötluse vahendid, mis suudavad analüüsida ja parandada mahukaid tekste märkimisväärse usaldusväärsusega~\cite{brown_language_2020, openai_gpt4_2023, anthropic_claude_2024}. Kui veel mõni aasta tagasi olid sellised vahendid valdavalt ingliskeelsed, siis tänapäevased mitmekeelsed mudelid, nagu Claude 4.7 Opus, GPT-5.5 ja Gemini 2.x, suudavad mõistliku kvaliteediga töödelda ka eesti keelt, mida kinnitavad TartuNLP avalikud võrdlused (vt baromeeter~\cite{tartunlp_baromeeter_2025}). See loob võimaluse rakendada LLM-e tudengite toetamiseks erialadel, kus eestikeelne akadeemiline kirjutamine on traditsiooniliselt olnud keeruline.

Akadeemilises kontekstis on LLM-ide kasutamine samas eetiliselt tundlik teema. Tartu Ülikooli generatiivse tehisaru kasutamise juhend~\cite{ut_genAI_2024} eristab selgelt teksti \emph{loomist} (mis ei ole lõputöös tudengi enda panusena lubatud) ja teksti \emph{tagasisidestamist} (mis on lubatud, kui tudeng saadud soovitusi ise kriitiliselt hindab ja rakendab). Käesolevas töös käsitletav süsteem kuulub üheselt teise kategooriasse: see ei kirjuta lõputööd tudengi eest ega asenda juhendajat, vaid annab tudengi enda kirjutatud tekstile soovituslikku tagasisidet. See eristus on teema akadeemilise vastuvõetavuse jaoks kriitilise tähtsusega ning seda tuleb selgelt rõhutada nii süsteemi kasutajaliideses kui ka tagasiside tekstis endas.

\subsection*{Töö eesmärk ja ulatus}

Käesoleva bakalaureusetöö \textbf{eesmärk} on \emph{kavandada ja realiseerida} LLM-il põhinev prototüüp, mis annab eestikeelsele informaatika lõputöö tekstile automaatset tagasisidet struktuuri, akadeemilise stiili, terminoloogia järjepidevuse ja viitamisvajaduse kohta, ning \emph{hinnata} selle väljundi kvaliteeti nii kvalitatiivse pilootuuringu kui ka sünteetilisel testkogul tehtud empiirilise hindamisega.

\textbf{Töö ulatus.} Käesoleva töö raames tehti:
\begin{enumerate}
    \item tagasisidekategooriate operatsionaalne määratlus;
    \item kahe promptimisstrateegia (üldine ja struktureeritud) disain;
    \item veebipõhise prototüübi realiseerimine, sealhulgas demo-režiim, mis võimaldab süsteemi proovida ilma välise LLM-i API juurdepääsuta;
    \item kvalitatiivne pilootuuring väikesel tekstivalimil prompti disaini iteratsiooniks ja mudeli väljundi tüüpiliste mustrite tuvastamiseks;
    \item 20-katkendilise sünteetilise testkogu ja kuldstandardi koostamine (autori juhendamisel LLM-iga; vt allpool generatiivse tehisaru deklaratsioon);
    \item täismahuline empiiriline hindamine 80 päris API-päringu põhjal (kaks mudelit × kaks prompti × 20 katkendit) koos täpsuse, saagise ja F\textsubscript{1}-skoori arvutamisega iga (mudel × promptitüüp × kategooria) kombinatsiooni kohta.
\end{enumerate}

\textbf{Töö ulatusest on teadlikult välja jäetud}: päris üliõpilastöödel põhinev testkogu (käesolev testkogu on sünteetiline) ning sõltumatu teise annoteerija kaasamine kuldstandardi koostamisse (käesoleva töö kuldstandard on ühe annoteerija ehk autori koostatud). Nende sammude jätkukava on esitatud peatükis~\ref{hindamine:edasine}. Selline ulatusepiirang on motiveeritud sellega, et päris üliõpilastööde põhjal kuldstandardi koostamine koos sõltumatu teise annoteerijaga nõuab eraldi tööd, mis kuulub edasisse uurimisse (näiteks magistriõppes või avaldatavas artiklis).

\subsection*{Uurimisküsimused}

Käesolev töö käsitleb järgmisi \textbf{uurimisküsimusi}, millele antakse vastused nii kvalitatiivse pilootuuringu (peatükk~\ref{hindamine:arutelu}) kui ka sünteetilisel testkogul tehtud kvantitatiivse empiirilise hindamise põhjal (peatükk~\ref{hindamine:synteetiline}); sama metoodika kohaldamine päris üliõpilastöödele kuulub jätkukavasse (peatükk~\ref{hindamine:edasine}):
\begin{enumerate}
    \item Millise iseloomuga tagasisidet annab tippkvaliteediga LLM eestikeelse informaatika lõputöö tekstile, kui talle antakse struktureeritud kategoorianõuetega prompt?
    \item Millistes tagasisidekategooriates on mudeli väljund kõige usaldusväärsem ning millistes kõige problemaatilisem, ja miks?
    \item Kas struktureeritud promptimine annab üldisest promptist süsteemsemat ja kasutuskõlblikumat tagasisidet?
\end{enumerate}

Töö keskne \textbf{hüpotees} -- mida uurimisküsimus~3 operatsionaliseerib -- on, et struktureeritud kategooria- ja skeemipõhine promptimine annab eestikeelse informaatika lõputöö tekstile tagasiside andmisel üldisest viibast süsteemsemat ja kasutuskõlblikumat väljundit (vt täpsem sõnastus peatükis~\ref{taust:promptimine}).

Pilootuuringu tähelepanekud ja empiirilised tulemused, mis nendele küsimustele vastavad, on esitatud peatükis~\ref{hindamine}.

\subsection*{Praktiline panus ja piirangud}

Töö praktiline tulemus on lokaalselt käivitatav veebiprototüüp, kus kasutaja saab valida lõputöö peatükitüübi (sissejuhatus, taust, metoodika, tulemused, kokkuvõte), sisestada enda kirjutatud teksti ning saada struktureeritud tagasisidet neljas eelnimetatud kategoorias. Prototüüp on kavandatud nii, et seda saab proovida \emph{demo-režiimis} salvestatud näidisvastustega ilma välise LLM-i API juurdepääsuta, mis muudab töö praktilise tulemuse reprodutseeritavaks minimaalsete eeldustega.

\textbf{Töö ei käsitle} LLM-i kasutamist lõputöö \emph{loomiseks}. Süsteem on mõeldud tudengi enda kirjutatud tekstile tagasiside saamiseks ega püüa asendada juhendaja sisulist tagasisidet. Mudeli soovitused on kasutaja jaoks vaid lähtepunkt -- iga soovituse vastuvõtmine, ümbersõnastamine või tagasilükkamine on tudengi enda otsustada.

\subsection*{Töö ülesehitus}

Peatükis~\ref{taust} esitatakse ülevaade suurte keelemudelite arhitektuurist, eestikeelsest LLM-tugist ja akadeemilise teksti automaatsest tagasisidestamisest, samuti kirjeldatakse seotud töid. Peatükk~\ref{metoodika} määratleb tagasisidekategooriad ning kirjeldab prompti disaini põhimõtteid ja hindamismetoodikat. Peatükis~\ref{prototuup} kirjeldatakse prototüübi arhitektuuri, tehnoloogilist lahendust ja kasutajaliidese disaini. Peatükk~\ref{hindamine} esitab pilootuuringu kvalitatiivsed tähelepanekud, sünteetilisel testkogul tehtud empiirilise hindamise tulemused ning päris üliõpilastöödel põhineva jätkuhindamise kava. Peatükis~\ref{kokkuvõte} antakse kokkuvõte töö panusest ning visandatakse võimalused edasiseks uurimiseks. Lisas~I on esitatud kasutatud sünteetilise testkogu kirjeldus ja kavandatud päris testkogu disain, lisas~II struktureeritud prompti täielik tekst, lisas~III lähtekoodi ja andmestiku struktuur ning lisas~IV prototüübi kasutusjuhend (sealhulgas demo-režiimi käivitamine ilma API-võtmeta).

\subsection*{Generatiivse tehisaru kasutamise deklaratsioon}

Käesoleva töö koostamisel on generatiivset tehisaru (peamiselt Anthropici mudelit Claude 4.7 Opus) kasutatud järgmisel kolmel viisil, mille autor selgesõnaliselt deklareerib:

\begin{enumerate}
    \item \textbf{Tekstilise abivahendina} -- loetavuse parandamiseks, sõnastuse poleerimiseks ja struktuurinõustamiseks. Kõik sisulised mõtted, töö ulatuse valikud, metoodilised valikud ja järeldused on autori enda otsused. Mudeli ettepanekud vaadati iga kord kriitiliselt üle ning sõnastati vajaduse korral ümber või jäeti kõrvale.
    \item \textbf{Hindamise testkogu ja kuldstandardi koostamise abivahendina} -- alapeatükis~\ref{hindamine:synteetiline} esitatud empiirilise hindamise jaoks on 20 testkatkendit (igaüks ligikaudu 100--160 sõna) ja 42 kuldstandardi annotatsiooni koostatud autori juhendamisel mudeliga Claude 4.7 Opus \emph{interaktiivse vestlusliidese kaudu (Claude Code)}, mitte programmiliste API-päringute kaudu. See ühekordne genereerimisprotsess on metoodiliselt eraldatud alapeatükis~\ref{hindamine:synteetiline} kirjeldatud hindamispäringutest, mis kasutavad Anthropici ja OpenAI tasulisi API-sid skriptipõhiselt; testkogu ja kuldstandardi koostamine ei toimunud nende API-de kaudu. Tegu ei ole anonüümitud üliõpilastööde katkenditega, vaid sünteetiliste tekstidega, mis on koostatud nii, et katkendid sisaldaksid ette nähtud tüüpilisi probleeme tagasisidekategooriates. See piirang on selgesõnaliselt välja toodud alapeatükis~\ref{hindamine:synteetiline} ning piirangute jaotises~\ref{hindamine:piirangud}. Päris üliõpilastööde katkendite kasutamine kuulub jätkukavasse (vt alapeatükk~\ref{hindamine:edasine}).
    \item \textbf{Hinnatavate mudelivastuste tegelike API-päringute kaudu} -- alapeatükis~\ref{hindamine:synteetiline} esitatud F\textsubscript{1}-skoorid \emph{on tegelikud empiirilised mõõtmised}: iga (katkend × promptitüüp × mudel) kombinatsiooni jaoks tehti tegelik API-päring Anthropici mudelile Claude 4.7 Opus või OpenAI mudelile GPT-5.5 (kokku 80 päringut), salvestati toorvastus repositooriumi kausta \texttt{statistika/andmed/paris\_vastused/} ning skooriti tekstivahemiku tasemel kattuvuse järgi kuldstandardi vastu. \emph{Need numbrid kajastavad tegelikku mudelite sooritust} antud sünteetilisel testkogul, kuid nende üldistatavus päris üliõpilastöödele on piiratud, kuna testkogu ise on sünteetiline (vt punkt 2).
\end{enumerate}

Selline kasutus on otseses kooskõlas Tartu Ülikooli generatiivse tehisaru juhendiga~\cite{ut_genAI_2024}, mis lubab generatiivse tehisaru kasutamist õppetöös \emph{tingimusel, et selline kasutus on selgesõnaliselt deklareeritud}. Reprodutseerimiseks vajalikud skriptid, salvestatud API-vastused ja andmestik on saadaval töö repositooriumis (\url{https://github.com/UkuRenekKronbergs/TI_bakalaureusetoo}) kaustas \texttt{statistika/}. Lugeja saab seal kontrollida, kuidas iga arv on saadud (failid \texttt{andmestik\_synth.py}, \texttt{paris\_hindamine.py} ning vastuste JSON-id alamkaustas \texttt{paris\_vastused/}).
